%% ma: L12-B01-0022
%% lop: 12
%% chuong: 1
%% bai: 1
%% dang: 12.1.9
%% loai: TN4
%% muc: TH
%% dap_an: "A"
%% nguon: "(NBV)-ĐỀ SỐ 12-MỖI NGÀY 1 ĐỀ THI 2025.pdf (D12, MỖI NGÀY 1 ĐỀ - Đề số 12), Phần I Câu 3"
%% kien_thuc: "Bài 1 (cực trị của hàm số, dấu của đạo hàm)"
%% trang_thai: chinh-thuc
%% ghi_chu: "Đáp án gốc A đúng."
\begin{ex}
Cho hàm số $y=f(x)$ có đạo hàm $f'(x)=(x+2)(x-3)$, $\forall x\in\mathbb R$. Điểm cực đại của hàm số đã cho là
\choice{\True $x=-2$}{$x=2$}{$x=3$}{$x=-3$}
\loigiai{$f'(x)=0\Leftrightarrow x=-2$ hoặc $x=3$. Do $f'(x)$ là tam thức bậc hai có hệ số cao nhất dương nên $f'(x)>0$ khi $x<-2$ hoặc $x>3$, $f'(x)<0$ khi $-2<x<3$.

Đạo hàm đổi dấu từ dương sang âm khi qua $x=-2$ nên đó là điểm cực đại (còn $x=3$ là điểm cực tiểu). Chọn A.}
\end{ex}
